% -*- coding=utf-8 -*-
%\newcommand\clcNumber{O175.3}      	% 分类号
\newcommand\studentId{你的学号}     		% 学号
\newcommand\thesisTitleCN{基于粒子群算法的第一和第二特征值及其单位特征向量的求解}   		% 中文题目
\newcommand\thesisTitleEN{Computation of the First and Second Eigenvalues and Their Corresponding Unit Eigenvectors Based on Particle Swarm Optimization}   		% 英文题目
\newcommand\thesisAuthor{你的名字}   		% 作者
\newcommand\supervisor{王俊}        		% 指导教师
\newcommand\majorclass{信息与计算科学}    % 一级学科名称
\newcommand\institute{理学院}            % 二级学科名称
\newcommand\thesisDateOne{2026 年 6 月}    % 提交日期：封面是2026年6月，但论文题目页是二零二六年六月
\newcommand\thesisDateTwo{二零二六年六月}    % 提交日期：封面是2026年6月，但论文题目页是二零二六年六月

%============================================如非必要，以下内容请勿动
{   
    \heiti                                                                  % 设置封面主要字体为黑体
    \setlength\parindent{0em}                                               % 设置首行缩进为0
    \ThisTileWallPaper{\paperwidth}{\paperheight}{figure/cover1.pdf}         % 设置封面背景
%    \vspace*{0.2cm}                                                         % 设置页面内容上间距
%    {
%        \zihao{-3}
%        \begin{tabular}{l@{}c p{3cm} l@{}c}     % @{}去除边框距离，更符合学校模板的下划线紧贴描述文字
%            {分类号:}\quad  & \underline{\makebox[3cm][c]{}} \\
%        \end{tabular}
%        \renewcommand{\arraystretch}{1}
%    }
    \vspace*{9.5cm}
%    \begin{center}
%        \zihao{-3}  
%        题目：\underline{\makebox[9cm][c]{\songti \thesisTitle}}
%    \end{center}
%% 如果你的标题太长，想分两行，可以把上面这四行代码换成下面这段注释了的代码。
    % \begin{center}
    %     {\zihao{-3}  题目：}   \underline{\makebox[22em][c]{\zihao{-3}\kaishu LHC能区pp碰撞下$J/\psi$产生的探究}} \par
    %     {\zihao{-3} \qquad\quad }   \underline{\makebox[22em][c]{\zihao{-3}\kaishu LHC能区pp碰撞下$J/\psi$产生的探究}} \par
    % \end{center}
    \vspace*{4.5cm}
    \renewcommand{\baselinestretch}{2}\selectfont                           % 设置声明的行间距
    {
        \begin{center}
            {\zihao{-3}  学\;\;\;\;\;\;\;\;院}   \underline{\makebox[15em][c]{\zihao{-3}\kaishu\institute}} \par
            {\zihao{-3}  专\;\;\;\;\;\;\;\;业}         \underline{\makebox[15em][c]{\zihao{-3}\kaishu\majorclass}} \par
            {\zihao{-3}  学生姓名} \underline{\makebox[15em][c]{\zihao{-3}\kaishu\thesisAuthor}} \par
            {\zihao{-3}  班级学号}         \underline{\makebox[15em][c]{\zihao{-3}\kaishu\studentId}} \par
            {\zihao{-3}  指导教师} \underline{\makebox[15em][c]{\zihao{-3}\kaishu\supervisor}} \par
        \end{center}
    }
    \vspace{1.2cm}
    \begin{center}
        \zihao{-3} \bfseries 
        \makebox[5cm][c]{\songti \thesisDateOne}
    \end{center}
}



\newpage
%============================================如非必要，以下内容请勿动
{   
    \heiti                                                                  % 设置封面主要字体为黑体
    \setlength\parindent{0em}                                               % 设置首行缩进为0
    \ThisTileWallPaper{\paperwidth}{\paperheight}{figure/cover2.pdf}         % 设置封面背景
    \vspace*{2.5cm}                                                         % 设置页面内容上间距
%    \vspace{2.5cm}
	\begin{center}
         % \zihao{-2}  \bfseries 
         % \makebox[9cm][c]{\songti \thesisTitleCN} %中文题目\underline{\makebox[9cm][c]{\songti \thesisTitleCN}}
		
		% \zihao{-2}  \bfseries 
		% \makebox[9cm][c]{\songti \thesisTitleEN} %英语题目\underline{\makebox[9cm][c]{\songti \thesisTitleCN}}
%%% 如果你的标题太长，想分两行，可以把上面这四行代码换成下面这段注释了的代码。
   \begin{center}
       {\zihao{-2}}    \makebox[22em][c]{\songti \zihao{-2}\bfseries 基于粒子群算法的第一和第二特征值及其单位}\par
       {\zihao{-2} \qquad }   \makebox[22em][c]{\songti \zihao{-2}\bfseries 特征向量的求解} \par
   \end{center}

%%% 如果你的标题太长，想分两行，可以把上面这四行代码换成下面这段注释了的代码。
   \begin{center}
       {\zihao{-2}  }   \makebox[22em][c]{\songti \zihao{-2}\bfseries Computation of the First and Second Eigenvalues and Their } \par
       {\zihao{-3} \qquad }   \makebox[22em][c]{\songti \zihao{-2}\bfseries Corresponding Unit Eigenvectors Based on} \par
      {\zihao{-3} \qquad \quad }   \makebox[22em][c]{\songti \zihao{-2}\bfseries Particle Swarm Optimization} \par
   \end{center}
   
% %%% 如果你的标题太长，想分两行，可以把上面这四行代码换成下面这段注释了的代码。
%    \begin{center}
%        {\zihao{-3}  题目：}   \underline{\makebox[22em][c]{\zihao{-3}\kaishu LHC能区pp碰撞下$J/\psi$产生的探究}} \par
%        {\zihao{-3} \qquad\quad }   \underline{\makebox[22em][c]{\zihao{-3}\kaishu LHC能区pp碰撞下$J/\psi$产生的探究}} \par
%    \end{center}

% %%% 如果你的标题太长，想分两行，可以把上面这四行代码换成下面这段注释了的代码。
%    \begin{center}
%        {\zihao{-3}  英文题目：}   \underline{\makebox[22em][c]{\zihao{-3}\kaishu LHC能区pp碰撞下$J/\psi$产生的探究}} \par
%        {\zihao{-3} \qquad\quad }   \underline{\makebox[22em][c]{\zihao{-3}\kaishu LHC能区pp碰撞下$J/\psi$产生的探究}} \par
%    \end{center}
    \end{center}
 %%%%%%%%%%%%%%%%%%%%%%%%%%%%%%%%%%%%%%%%%%%%%%%%%%%%%%%%%%%%%%%%%%%%%%%%%%%%%%%%%%%%%%%%%%%%%%%%%%%%%%%%%%%%%%%%%%%%%
    \vspace*{4.5cm}
    \renewcommand{\baselinestretch}{2}\selectfont                           % 设置声明的行间距
%    {
        \begin{center}
            {\zihao{-3}  学生姓名：} \makebox[1.5cm][l]{\zihao{-3}\kaishu\thesisAuthor} \par
            {\zihao{-3}  指导教师：} \makebox[1.5cm][l]{\zihao{-3}\kaishu\supervisor} \par
		\end{center}
		
		
		\vspace{1.2cm}
	    \begin{center}
		\zihao{-3} %\bfseries 
        \makebox[5cm][c]{\songti 江苏科技大学}\par
		
        \zihao{-3} %\bfseries 
        \makebox[5cm][c]{\songti \thesisDateTwo}\par
		\end{center}
%    }
}
%{   
%    \heiti                                                                  % 设置封面主要字体为黑体
%    \setlength\parindent{0em}                                               % 设置首行缩进为0
%    \ThisTileWallPaper{\paperwidth}{\paperheight}{figure/cover2.pdf}         % 设置封面背景
%    \vspace*{0.2cm}                                                         % 设置页面内容上间距
%    \vspace{4.5cm}
%	\begin{center}
%        \zihao{-3}  \bfseries 
%        \makebox[9cm][c]{\songti \thesisTitleCN} %中文题目\underline{\makebox[9cm][c]{\songti \thesisTitleCN}}
%    \end{center}
%%% 如果你的标题太长，想分两行，可以把上面这四行代码换成下面这段注释了的代码。
%    % \begin{center}
%    %     {\zihao{-3}  题目：}   \underline{\makebox[22em][c]{\zihao{-3}\kaishu LHC能区pp碰撞下$J/\psi$产生的探究}} \par
%    %     {\zihao{-3} \qquad\quad }   \underline{\makebox[22em][c]{\zihao{-3}\kaishu LHC能区pp碰撞下$J/\psi$产生的探究}} \par
%    % \end{center}
%    \vspace*{4.5cm}
%    \renewcommand{\baselinestretch}{2}\selectfont                           % 设置声明的行间距
%    {
%        \begin{center}
%            {\zihao{-3}  学生姓名：} \makebox[15em][l]{\zihao{-3}\kaishu\thesisAuthor} \par
%            {\zihao{-3}  指导教师：} \makebox[15em][l]{\zihao{-3}\kaishu\supervisor} \par
%        \end{center}
%    }
%    \vspace{1.2cm}
%    \begin{center}
%	    \zihao{-3} %\bfseries 
%        \makebox[5cm][c]{\songti 江苏科技大学} \\
%        \zihao{-3} %\bfseries 
%        \makebox[5cm][c]{\songti \thesisDate}
%    \end{center}
%}